\section{USAMO 2024}
        \begin{enumerate}
        	\item (\textit{USAMO 2024}) Find all integers $n \geq 3$ such that the following property holds: if we list the divisors of $n !$ in increasing order as $1=d_1<d_2<\cdots<d_k=n!$, then we have


 
 \[d_2-d_1 \leq d_3-d_2 \leq \cdots \leq d_k-d_{k-1}.\]


 



\textbf{Solution (Explanation of Video Solution)}

We can start by verifying that $n=3$ and $n=4$ work by listing out the factors of $3!$ and $4!$. We can also see that $n=5$ does not work because the terms $15, 20$, and $24$ are consecutive factors of $5!$. Also, $n=6$ does not work because the terms $6, 8$, and $9$ appear consecutively in the factors of $6!$.

 Note that if we have a prime number $p>n$ and an integer $k>p$ such that both $k$ and $k+1$ are factors of $n!$, then the condition cannot be satisfied.

 If $n\geq7$ is odd, then $(2)(\frac{n-1}{2})(n-1)=n^2-2n+1$ is a factor of $n!$. Also, $(n-2)(n)=n^2-2n$ is a factor of $n!$. Since $2n<n^2-2n$ for all $n\geq7$, we can use \href{https://artofproblemsolving.com/wiki/index.php?title=Bertrand\%27s\_Postulate}{Bertrand's Postulate} to show that there is at least one prime number $p$ such that $n<p<n^2-2n$. Since we have two consecutive factors of $n!$ and a prime number between the smaller of these factors and $n$, the condition will not be satisfied for all odd $n\geq7$.

 If $n\geq8$ is even, then $(2)(\frac{n-2}{2})(n-2)=n^2-4n+4$ is a factor of $n!$. Also, $(n-3)(n-1)=n^2-4n+3$ is a factor of $n!$. Since $2n<n^2-4n+3$ for all $n\geq8$, we can use Bertrand's Postulate again to show that there is at least one prime number $p$ such that $n<p<n^2-4n+3$. Since we have two consecutive factors of $n!$ and a prime number between the smaller of these factors and $n$, the condition will not be satisfied for all even $n\geq8$.

 Therefore, the only numbers that work are $n=3$ and $n=4$.

 As listed above, $n=3$ and $n=4$ work but $n=5$ and $n=6$ don't. Assume $n>6$, let $p$ be the smallest positive integer that doesn't divide $n!$. It is easy to see $p$ is the smallest prime greater than $n$. It is easy to see that both $p-1$ and $p+1$ are divisors of $n!$ Let $q$ be the smallest positive integer greater than $p$ that is 2 mod 6. We see that either $q=p+3$ or $q=p+1$ but $q$, $q+1$, $q+2$ must all be divisors of $n!$ giving a difference of 1 after a difference of 2. Therefore, all $n>6$ fail. ~mathophobia

 



\textbf{Solution 2}

We claim only $n = 3$ and $n = 4$ are the only two solutions. First, it is clear that both solutions work. 

 Next, we claim that $n < 5$. For $n \geq 5$, let $x$ be the smallest $x$ such that $x+1$ is not a factor of $n!$. Let the smallest factor larger than $x$ be $x+k$. 

 Now we consider $\frac{n!}{x-1}$, $\frac{n!}{x}$ and $\frac{n!}{x+k}$. Since $\frac{n!}{x-1} > \frac{n!}{x} > \frac{n!}{x+k}$, if $n$ were to satisfy the conditions, then $\frac{n!}{x-1}-\frac{n!}{x} \geq \frac{n!}{x} - \frac{n!}{x+k}$. However, note that this is not true for $x \geq 5$ and $k > 1$. 

 Note that the inequality is equivalent to showing $\frac{1}{x(x-1)} \geq \frac{k}{x(x+k)}$, which simplifies to $x+k \geq kx-k$, or $\frac{x}{x-2} \geq k \geq 2$. This implies $x \geq 2x-4$, $x \leq 4$, a contradiction, since the set of numbers $\{1, 2, 3, 4, 5\}$ are all factors of $n!$, and the value of $x$ must exist. Hence, no solutions for $n \geq 5$.

 


	\item (\textit{USAMO 2024}) Let $S_1, S_2, \ldots, S_{100}$ be finite sets of integers whose intersection is not empty. For each non-empty $T \subseteq\left\{S_1, S_2, \ldots, S_{100}\right\}$, the size of the intersection of the sets in $T$ is a multiple of the number of sets in $T$. What is the least possible number of elements that are in at least 50 sets?


 



\textbf{Solution 1}

Let's determine the smallest possible number of elements that appear in at least 50 of the sets.

 First, we establish some notation. For any subset $T \subseteq \{S_1, S_2, \ldots, S_{100}\}$, let $\cap T$ denote the intersection of all sets in $T$. By the problem statement, $|\cap T|$ is divisible by $|T|$.

 We'll use a binary encoding approach. For each element $x$ in our universe, define its "signature" as a binary string of length 100, where the $i$-th bit is 1 if $x \in S_i$ and 0 otherwise.

 For any signature $v$ with exactly $k$ ones, corresponding elements must appear in exactly $k$ sets. The problem condition requires that for any subset of sets corresponding to positions where 1's appear in some signature $u$, the number of elements having signature $v$ that contains all 1's from $u$ must be divisible by $|u|$.

 Let's denote by $f(v)$ the number of elements with signature $v$. The problem condition translates to: for any signature $u$, $|u|$ divides $\sum_{v \supseteq u} f(v)$.

 Key claim: For any signature $u$ with $|u| = 50$, we need $f(u) \geq 50$.

 Proof: By the divisibility condition, $\sum_{v \supseteq u} f(v)$ must be divisible by 50. For any $v$ strictly containing $u$, if we sum over all such $u$ of size 50, each $f(v)$ with $|v| > 50$ gets counted multiple times. To satisfy the divisibility condition for each individual $u$, we need $f(u) \geq 50$ for each signature $u$ with exactly 50 ones.

 Since there are $\binom{100}{50}$ different ways to choose 50 positions for ones in the signature, and each such signature must correspond to at least 50 elements, the minimum number of elements appearing in at least 50 sets is $50 \binom{100}{50}$.

 Therefore, the answer is $\boxed{50 \binom{100}{50}}$.

 


	\item (\textit{USAMO 2024}) Let $m$ be a positive integer. A triangulation of a polygon is $m$-balanced if its triangles can be colored with $m$ colors in such a way that the sum of the areas of all triangles of the same color is the same for each of the $m$ colors. Find all positive integers $n$ for which there exists an $m$-balanced triangulation of a regular $n$-gon.


 Note: A triangulation of a convex polygon $\mathcal{P}$ with $n \geq 3$ sides is any partitioning of $\mathcal{P}$ into $n-2$ triangles by $n-3$ diagonals of $\mathcal{P}$ that do not intersect in the polygon's interior.


 



\textbf{Solution 1}

The answer is if and only if $m$ is a proper divisor of $n$.

 We represent the vertices of the regular $n$-gon using complex numbers. Let $\omega = \exp(2\pi i/n)$ be a primitive $n$-th root of unity, and label the vertices as $1, \omega, \omega^2, \ldots, \omega^{n-1}$.

 Key Lemma:

 The triangle with vertices $\omega^k, \omega^{k+a}, \omega^{k+b}$ has signed area:

 
 \[T(a, b) := \frac{(\omega^a - 1)(\omega^b - 1)(\omega^{-b} - \omega^{-a})}{4i}\]

 Proof of Lemma:

 We can rotate by $\omega^{-k}$ to assume without loss of generality that $k = 0$. Then we apply the complex shoelace formula to the triangle with vertices $1, \omega^a, \omega^b$:

 
 \[\frac{i}{4} \det \begin{bmatrix} 1 & 1 & 1 \\ \omega^a & \omega^{-a} & 1 \\ \omega^b & \omega^{-b} & 1 \end{bmatrix} = \frac{i}{4} \det \begin{bmatrix} 0 & 0 & 1 \\ \omega^a - 1 & \omega^{-a} - 1 & 1 \\ \omega^b - 1 & \omega^{-b} - 1 & 1 \end{bmatrix}\]

 which simplifies to the expression for $T(a, b)$ above.

 Construction (Sufficiency):

 To construct a valid triangulation and coloring when $m$ divides $n$, we draw all the diagonals from vertex 1, and then color the resulting triangles with the $m$ colors in cyclic order.

 For example, when $n = 9$ and $m = 3$, we get a coloring with red, green, blue as shown in the image, repeating cyclically.

 To prove this works, we fix a residue $r \mod m$ corresponding to one of the colors. Then the total area of triangles with this color is:

 
 \[\sum_{\substack{0 \leq j < n \\ j \equiv r \mod m}} \text{Area}(\omega^0, \omega^j, \omega^{j+1}) = \sum_{\substack{0 \leq j < n \\ j \equiv r \mod m}} T(j, j+1)\]

 After algebraic manipulation:

 
 \[= \frac{\omega - 1}{4i} \left(\frac{n}{m}(1 + \omega^{-1}) + \sum_{\substack{0 \leq j < n \\ j \equiv r \mod m}} (\omega^{-1-j} - \omega^j)\right)\]

 When $m$ divides $n$, we can show that the inner sum vanishes, and the total area of each color equals:

 
 \[\sum_{\substack{0 \leq j < n \\ j \equiv r \mod m}} \text{Area}(\omega^0, \omega^j, \omega^{j+1}) = \frac{n}{m} \cdot \frac{\omega - \omega^{-1}}{4i}\]

 Since the right-hand side does not depend on the residue $r$, this shows all colors have equal area.

 Necessity:

 To prove that $m$ must divide $n$, we note that if there were a valid triangulation and coloring, the total area of each color would equal:

 
 \[S := \frac{n}{m} \cdot \frac{\omega - \omega^{-1}}{4i}\]

 We then prove the following claim:

 The number $4i \cdot S$ is not an algebraic integer when $m \nmid n$.

 This can be shown using the fact that $K = \mathbb{Q}(\omega)$ is a number field with ring of integers $\mathcal{O}_K = \mathbb{Z}[\omega]$. We demonstrate that:

 
 \[\omega \cdot 4i \cdot S = \frac{n}{m}\omega^2 - \frac{n}{m}\]

 is not an algebraic integer when $m$ doesn't divide $n$ because $\frac{n}{m} \not\in \mathbb{Z}$.

 However, for any triangle in our construction, $4i \cdot T(a,b)$ is always an algebraic integer. Since a finite sum of algebraic integers is also an algebraic integer, the sum of expressions of the form $4i \cdot T(a,b)$ can never equal $4i \cdot S$ unless $m$ divides $n$.

 We can also show this directly by noting that $T(a,b)$ is always divisible by $\frac{\omega - \omega^{-1}}{4i}$, which means we need $\frac{n}{m}$ to be an algebraic integer. A rational number is an algebraic integer if and only if it's a rational integer, so $\frac{n}{m}$ is an algebraic integer exactly when $m$ divides $n$.

 Therefore, a valid triangulation and coloring exists if and only if $m$ is a proper divisor of $n$.

 


	\item (\textit{USAMO 2024}) Let $m$ and $n$ be positive integers. A circular necklace contains $m n$ beads, each either red or blue. It turned out that no matter how the necklace was cut into $m$ blocks of $n$ consecutive beads, each block had a distinct number of red beads. Determine, with proof, all possible values of the ordered pair $(m, n)$.


 



\textbf{Solution 1}

We need to determine all possible positive integer pairs $(m, n)$ such that there exists a circular necklace of $mn$ beads, each colored red or blue, satisfying the following condition:

 No matter how the necklace is cut into $m$ blocks of $n$ consecutive beads, each block has a distinct number of red beads.

 Necessary Condition:

 1. Maximum Possible Distinct Counts:  In a block of $n$ beads, the number of red beads can range from $0$ to $n$.  Therefore, there are $n + 1$ possible distinct counts of red beads in a block.  Since we have $m$ blocks, the maximum number of distinct counts must be at least $m$.  Thus, we must have:  
 \[m \leq n + 1\]

 Sufficient Construction:

 We will show that for all positive integers $m$ and $n$ satisfying $m \leq n + 1$, such a necklace exists.

 1. Construct Blocks:  Create $m$ blocks, each containing $n$ beads. Assign to each block a unique number of red beads, ranging from $0$ to $m - 1$.

 2. Design the Necklace:  Arrange these $m$ blocks in a fixed order to form the necklace. Since the necklace is circular, cutting it at different points results in cyclic permutations of the blocks.

 3. Verification:  In any cut, the sequence of blocks (and thus the counts of red beads) is a cyclic shift of the original sequence. Therefore, in each partition, the blocks will have distinct numbers of red beads.

 Example:

 Let's construct a necklace for $m = 3$ and $n = 2$:

 Blocks:  Block 1: $0$ red beads (BB)  Block 2: $1$ red bead (RB)  Block 3: $2$ red beads (RR)

 Necklace Arrangement:  Place the blocks in order: BB-RB-RR.

 Verification:  Any cut of the necklace into $3$ blocks of $2$ beads will have blocks with red bead counts of $0$, $1$, and $2$.

 Conclusion:

 All ordered pairs $(m, n)$ where $m \leq n + 1$ satisfy the condition.  Therefore, the possible values of $(m, n)$ are all positive integers such that $m \leq n + 1$.

 Final Answer:

 Exactly all positive integers $m$ and $n$ with $m \leq n + 1$; these are all possible ordered pairs $(m, n)$.

 


	\item (\textit{USAMO 2024}) Point $D$ is selected inside acute triangle $A B C$ so that $\angle D A C=$ $\angle A C B$ and $\angle B D C=90^{\circ}+\angle B A C$. Point $E$ is chosen on ray $B D$ so that $A E=E C$. Let $M$ be the midpoint of $B C$.
Show that line $A B$ is tangent to the circumcircle of triangle $B E M$.


 



\textbf{Solution 1}

Let $\angle DBT = \alpha$ and $\angle BEM = \beta$.  Extend AD intersects BC at point T, then TC = TA, TE is perpendicular to AC

 Thus, AB is the tangent of the circle BEM

 Then the question is equivalent as  the $\angle ABT$ is the auxillary angle of $\angle BEM$.

 \textit{ontinued}

 



\textbf{Solution 2}

We notice that $\angle DAC = \angle ACB$, which implies symmetry about the perpendicular bisector of $\overline{AC}$. Therefore, we can extend $\overline{AD}$ to intersect the circumcircle at $Q$, and by symmetry about the perpendicular bisector, $ABQC$ is an isoceles trapezoid. Now, we realize that the midpoint condition is annoying, which really suggests that to simplify things, we take a homothety at $B$ sending $M$ to $C$, and we just need to prove that the circumcircle of $\triangle BE'C$ is tangent to line $\overline{AB}$, where $E'$ is the resultant point from a homothety at $E$. However, funnily enough, $E$ lies on the perpendicular bisector of $\overline{AC}$ because $\overline{AE}=\overline{EC}$, and the perpendicular bisector gets sent to an altitude from $Q$ under homothety, which means $E'$ is simply the intersection of $\overline{BD}$ and the altitude from $Q$ onto $\overline{BC}$! Call that point $P$. We can now ignore points $E$ and $M$. Note that we have not used the $\angle BDC$ condition yet, which means we should probably try angle chasing now. Call $\angle BAC=a, \angle ABC = b$. By angle condition, $\angle BDC = 90+a, \angle CDP = 90-a$. Now, $\angle QCA = \angle BAC = a$, so $\angle PQC = 90 - \angle ACQ = 90-a$. A miracle occurs! $\angle CDP = \angle ACQ$, so quadrilateral $QDPC$ is cyclic! $\angle ABC = \angle AQC = b$ from isoceles trapezoid, and from cyclicity $\angle DPC = 180-b = \angle BPC$. Now we have $\angle ABC = b$ and $\angle BPC = 180-b$, which looks a lot like tangency condition. We realize that it simply is the tangency condition if we extend line $\overline{AB}$ to get the supplementary angle, so we are done.

 


	\item (\textit{USAMO 2024}) Let $n>2$ be an integer and let $\ell \in\{1,2, \ldots, n\}$. A collection $A_1, \ldots, A_k$ of (not necessarily distinct) subsets of $\{1,2, \ldots, n\}$ is called $\ell$-large if $\left|A_i\right| \geq \ell$ for all $1 \leq i \leq k$. Find, in terms of $n$ and $\ell$, the largest real number $c$ such that the inequality

 \[\sum_{i=1}^k \sum_{j=1}^k x_i x_j \frac{\left|A_i \cap A_j\right|^2}{\left|A_i\right| \cdot\left|A_j\right|} \geq c\left(\sum_{i=1}^k x_i\right)^2\]
holds for all positive integers $k$, all nonnegative real numbers $x_1, \ldots, x_k$, and all $\ell$-large collections $A_1, \ldots, A_k$ of subsets of $\{1,2, \ldots, n\}$.


 Note: For a finite set $S,|S|$ denotes the number of elements in $S$.


 



\textbf{Solution}

The largest real number $c$ is $\frac{n - 2\ell + \ell^2}{n(n-1)}$.

 \textbf{Proof} For each $1 \leq i \leq k$, let $\mathbf{1}_{A_i} : \{1, 2, \ldots, n\} \to \{0, 1\}$ denote the indicator function of the set $A_i$. We can rewrite the intersection size in the numerator as:

 
 \[|A_i \cap A_j|^2 = \left( \sum_{m=1}^n \mathbf{1}_{A_i}(m) \mathbf{1}_{A_j}(m) \right)^2 = \sum_{m=1}^n \sum_{m'=1}^n \mathbf{1}_{A_i}(m) \mathbf{1}_{A_i}(m') \mathbf{1}_{A_j}(m) \mathbf{1}_{A_j}(m').\]

 Substitute this into the original double sum and rearrange the order of summation:
 \[\sum_{i=1}^k \sum_{j=1}^k x_i x_j \frac{|A_i \cap A_j|^2}{|A_i| \cdot |A_j|} = \sum_{m=1}^n \sum_{m'=1}^n \left( \sum_{i=1}^k \frac{x_i}{|A_i|} \mathbf{1}_{A_i}(m) \mathbf{1}_{A_i}(m') \right)^2.\]

 Define a symmetric $n \times n$ matrix $W$ where the entry in the $m$-th row and $m'$-th column is:
 \[W_{m, m'} = \sum_{i=1}^k \frac{x_i}{|A_i|} \mathbf{1}_{A_i}(m) \mathbf{1}_{A_i}(m').\]

 The expression we wish to bound is the sum of the squares of all entries in $W$, which we can split into diagonal and off-diagonal sums:
 \[\sum_{m, m'} W_{m, m'}^2 = \sum_{m=1}^n W_{m,m}^2 + \sum_{m \neq m'} W_{m, m'}^2.\]

 By the Cauchy-Schwarz inequality, the sum of squares of $N$ elements is bounded below by the square of their sum divided by $N$. Applying this to the $n$ diagonal elements and the $n(n-1)$ off-diagonal elements separately yields:
 \[\sum_{m=1}^n W_{m,m}^2 \geq \frac{1}{n} \left( \sum_{m=1}^n W_{m,m} \right)^2\]
 \[\sum_{m \neq m'} W_{m,m'}^2 \geq \frac{1}{n(n-1)} \left( \sum_{m \neq m'} W_{m,m'} \right)^2\]

 Next, we evaluate the sums of these entries. The sum of the diagonal elements (the trace) is:
 \[\sum_{m=1}^n W_{m,m} = \sum_{i=1}^k \frac{x_i}{|A_i|} \sum_{m=1}^n \mathbf{1}_{A_i}(m) = \sum_{i=1}^k \frac{x_i}{|A_i|} |A_i| = \sum_{i=1}^k x_i.\]

 The sum of all entries in the matrix is:
 \[\sum_{m, m'} W_{m, m'} = \sum_{i=1}^k \frac{x_i}{|A_i|} \left( \sum_{m=1}^n \mathbf{1}_{A_i}(m) \right)^2 = \sum_{i=1}^k x_i |A_i|.\]

 Subtracting the trace from the total sum gives the sum of the off-diagonal elements:
 \[\sum_{m \neq m'} W_{m, m'} = \sum_{i=1}^k x_i |A_i| - \sum_{i=1}^k x_i = \sum_{i=1}^k x_i (|A_i| - 1).\]

 Because the collection is $\ell$-large, we know $|A_i| \geq \ell \geq 1$. Given $x_i \geq 0$, we have:
 \[\sum_{m \neq m'} W_{m, m'} \geq \sum_{i=1}^k x_i (\ell - 1) = (\ell - 1) \sum_{i=1}^k x_i \geq 0.\]

 Let $S = \sum_{i=1}^k x_i$. Substituting these into our Cauchy-Schwarz bounds:
 \[\sum_{m, m'} W_{m, m'}^2 \geq \frac{1}{n} S^2 + \frac{1}{n(n-1)} (\ell - 1)^2 S^2 = \left( \frac{n - 1 + (\ell - 1)^2}{n(n-1)} \right) S^2 = \frac{n - 2\ell + \ell^2}{n(n-1)} S^2.\]

 This establishes that $c = \frac{n - 2\ell + \ell^2}{n(n-1)}$ is a valid lower bound. 

 To show this bound is achievable, consider the collection of all $k = \binom{n}{\ell}$ distinct subsets of size $\ell$, with $x_i = 1$ for all $i$. By symmetry, all $n$ diagonal entries of $W$ are identical, and all $n(n-1)$ off-diagonal entries are identical. Thus, the equality condition holds for both applications of Cauchy-Schwarz, confirming the maximum constant is exactly $\frac{n - 2\ell + \ell^2}{n(n-1)}$.

 


\end{enumerate}